% Original self-contained variational bridge.
\subsection*{Worked calculation: Gaussian upper bound for an anharmonic oscillator}
The helium calculation above illustrates screening, but its
Coulomb integrals can obscure the variational principle.
A one-dimensional example lets us calculate every integral
and every derivative explicitly before handling more elaborate atoms.

Choose dimensionless units $\hbar=m=\omega=1$ and
\begin{equation}
\op H=-\frac12\frac{\dd^2}{\dd x^2}
+\frac{x^2}{2}+\lambda x^4,\qquad\lambda\geq0.
\end{equation}
Take a normalized Gaussian trial state parametrized by
its positive width $\sigma$,
\begin{equation}
\psi_\sigma(x)=\frac{1}{\pi^{1/4}\sqrt\sigma}
 \exp\left(-\frac{x^2}{2\sigma^2}\right),
\qquad
|\psi_\sigma(x)|^2
=\frac{e^{-x^2/\sigma^2}}{\sqrt\pi\,\sigma}.
\end{equation}
Check normalization by setting $u=x/\sigma$:
\begin{equation}
\int_{-\infty}^{\infty}|\psi_\sigma|^2\dd x
=\frac1{\sqrt\pi\,\sigma}
 \int_{-\infty}^{\infty}e^{-x^2/\sigma^2}\dd x
=\frac1{\sqrt\pi}\int_{-\infty}^{\infty}e^{-u^2}\dd u
=1.
\end{equation}
The Gaussian integral is $\int_{\mathbb R}e^{-u^2}\dd u
=\sqrt\pi$. Differentiate it with respect to a width parameter
to obtain
\begin{equation}
\langle x^2\rangle
=\frac1{\sqrt\pi\,\sigma}
 \int_{\mathbb R}x^2e^{-x^2/\sigma^2}\dd x
=\frac{\sigma^2}{2},\qquad
\langle x^4\rangle=\frac{3\sigma^4}{4}.
\end{equation}
For the kinetic term, first differentiate
$\psi'_\sigma=-x\psi_\sigma/\sigma^2$.
Integration by parts, with the boundary term vanishing, gives
\begin{align}
\left\langle-\frac12\frac{\dd^2}{\dd x^2}\right\rangle
&=\frac12\int_{\mathbb R}|\psi'_\sigma(x)|^2\dd x\notag\\
&=\frac{1}{2\sigma^4}\int_{\mathbb R}
 x^2|\psi_\sigma(x)|^2\dd x
=\frac{1}{4\sigma^2}.
\end{align}
Combining these separate contributions,
\begin{equation}
\boxed{E_{\mathrm{trial}}(\sigma)
=\frac1{4\sigma^2}+\frac{\sigma^2}4
+\frac{3\lambda\sigma^4}{4}.}
\label{eq:worked-variational-energy}
\end{equation}
To find the optimum, calculate
\begin{equation}
\frac{\dd E_{\mathrm{trial}}}{\dd\sigma}
=-\frac1{2\sigma^3}+\frac{\sigma}2
+3\lambda\sigma^3.
\end{equation}
The stationarity condition can be multiplied by
$2\sigma^3$ (positive because $\sigma>0$):
\begin{equation}
\boxed{\sigma^4+6\lambda\sigma^6=1.}
\end{equation}
For $\lambda\geq0$, the left side is strictly increasing
in $\sigma>0$, ranging from zero to infinity. Thus it
has exactly one positive solution.
Furthermore,
\begin{equation}
\frac{\dd^2 E_{\mathrm{trial}}}{\dd\sigma^2}
=\frac3{2\sigma^4}+\frac12+9\lambda\sigma^2>0.
\end{equation}
The stationary point is therefore the unique global minimum.
The variational bound is
$E_0\leq E_{\mathrm{trial}}(\sigma_\ast)$.
At $\lambda=0$, the equation gives $\sigma_\ast=1$
and $E_{\mathrm{trial}}=1/2$, recovering the exact
harmonic-oscillator ground energy.

\paragraph{Why the inequality always points upward.}
Expand any \emph{normalized} trial ket in exact normalized
energy eigenstates:
$\ket\psi=\sum_n c_n\ket n$, so
$\sum_n|c_n|^2=1$. Then
\begin{equation}
\langle H\rangle-E_0
=\sum_n|c_n|^2(E_n-E_0)\geq0,
\end{equation}
because each term is nonnegative. If the trial ket were not
normalized, the correct variational object would be
$\mel\psi{\op H}\psi/\braket\psi\psi$.

\paragraph{Interpretation and experiment.}
Increasing $\lambda$ penalizes large $\langle x^4\rangle$,
so the optimal wavepacket becomes narrower. Its kinetic
energy then increases, and the width balances these effects.
The after-class
\href{https://github.com/KangqiaoLiu/advanced-quantum-mechanics-lecture-note/blob/main/companion/labs/IC-04-02/index.html}
{variational energy interactive investigation}
plots the same formula against an independent finite-basis
diagonalization, giving two variational upper bounds that
can be compared for convergence.
